% Context from chapters/approximation.tex; for mathematical reference, not compilation.
s $f \in \Ldeux([0,1]^d)$ of arbitrary dimension $d$ as
\eql{\label{eq-sobolev-norm-highdim}
	\norm{f}_{\text{Sob}(\al)}^2 = \sum_{m \in \ZZ^d} (2\pi \norm{m})^{2\al} |\dotp{f}{e_m}|^2,
}

The periodic Sobolev model
\eql{\label{eq-sobolev-model}
	\Th = \enscond{f \in \Ldeux([0,1]^d)}{ \norm{f}_2^2+\norm{f}_{\text{Sob}(\al)}^2\leq C }
}
includes the corresponding periodic $C^\alpha$ model~\eqref{eq-smooth-model} for integer $\alpha$, after adjusting the radius of the ball. For noninteger $\alpha$, H\"older regularity gives Sobolev regularity of every order $s<\alpha$, but need not give regularity of order $\alpha$ itself.

Figure \ref{fig-signal-model-sobolev} shows images with an increasing Sobolev norm for $\al=2$.

\myfigure{
\image{approximation}{.5}{signal-model-sobolev}
}{ 
	Images with increasing Sobolev norm.  
}{fig-signal-model-sobolev}


                                           
\subsection{Piecewise Regular Signals and Images}

  
\paragraph{Piecewise smooth signals.}

A piecewise smooth 1-D signal $f \in \Ldeux([0,1])$ is $\Cal$ on the intervals separated by at most $K$ interior singular points:
\eql{\label{eq-model-piece-smooth}
	\Th = \enscond{f \in \Ldeux([0,1])}{ \exists\,0=t_0<t_1<\cdots<t_{K+1}=1,\quad\max_{0\leq i\leq K}\norm{f|_{(t_i,t_{i+1})}}_{\Cal}\leq C }
}
Here $f|_{(t_i,t_{i+1})}$ denotes the restriction of $f$ to the open interval $(t_i,t_{i+1})$.

  
\paragraph{Piecewise smooth images.}

A piecewise smooth image is a function $f \in \Ldeux([0,1]^2)$ with $\Cal$ regularity away from at most $K$ rectifiable curves:
\eql{\label{eq-model-piece-smooth-2d}
	\Th = \enscond{f \in \Ldeux([0,1]^2)}{ \exists \, \Ga=\bigcup_{i=0}^{K-1}\ga_i, \;
	\norm{f}_{\Cal(\Ga^c)} \leq C_1 \qandq
	|\ga_i| \leq C_2
	 }	 
}
Here $|\ga_i|$ denotes curve length. The $C^\alpha$ bound applies separately within each smooth region, with uniform one-sided regularity up to its edges.

Segmentation methods such as the one proposed by Mumford and Shah \cite{mumford-shah} implicitly assume such a piecewise smooth image model.

                                           
\subsection{Bounded Variation Signals and Images}
\label{subsec-bv-model}

Bounded variation provides a more flexible model for signals with edges
\eql{\label{eq-model-tv}
	\Th=\enscond{f\in\Ldeux([0,1]^d)}{ \normi{f} \leq C_1 \qandq \normTV{f} \leq C_2 }.
}
For $d=1$ and $d=2$, this model includes the piecewise smooth signal and image models~\eqref{eq-model-piece-smooth} and~\eqref{eq-model-piece-smooth-2d} when $\alpha\geq1$, after adjusting the constants. Bounds on the derivatives, jump sizes, and total edge length then control the variation.

The total variation of a smooth function is 
\eq{
	\int \norm{\nabla f(x)} \d x
} 
where 
\eq{
	\nabla f(x) = \pa{ \frac{\partial f}{\partial x_i}  }_{i=0}^{d-1} \in \RR^d
} 
is the gradient at $x$. Total variation also extends to discontinuous images, including those with jumps across edges. The coarea formula expresses the total variation of a piecewise smooth image as an integral of the perimeters of its superlevel sets:
\eql{\label{eq-tv-coaera}
	\normTV{f}=\int_{-\infty}^{\infty}\operatorname{Per}(\{f>t\};\Omega)\,\d t,\qquad\Omega=(0,1)^d.
}	
Here the perimeter is relative to the image domain. For smooth functions in two dimensions it agrees, for almost every $t$, with the length of the level curve. For a bounded set $\Om\subset\RR^2$ with piecewise smooth boundary, taking total variation on all of $\RR^2$ gives
\eq{
	\normTV{1_\Om} = |\partial \Om|.
}
The model of bounded variation was introduced in image processing by Rudin, Osher and Fatemi \cite{rudin-tv}.

                                           
\subsection{Cartoon Images}
\label{subsec-cartoon-images}

The bounded-variation model~\eqref{eq-model-tv} controls total perimeter but does not require individual edges to be smooth. For images with smooth contours, incorporating that additional geometry can improve approximation and processing.

The $\Cal$ cartoon model consists of 2-D functions that are $\Cal$ away from at most $K$ smooth edge curves $\ga_i$:
\eql{\label{eq-cartoon-model}
	\Th = \enscond{f \in \Ldeux([0,1]^2)}{ \exists \, \Ga=\bigcup_{i=0}^{K-1}\ga_i, \;
	\norm{f}_{\Cal(\Ga^c)} \leq C_1 \qandq
	\norm{\ga_i}_{\Cal} \leq C_2
	 }
}
Each curve is parameterized by arc length on $[0,A_i]$, with a uniform bound on $A_i$ and its $C^\alpha$ norm.
Figure \ref{fig-image-model-cartoon} shows cartoon images with increasing total variation $\normTV{f}$.

\myfigure{
\image{approximation}{.5}{image-model-cartoon}
}{ 
	Cartoon images with increasing total variation. 
}{fig-image-model-cartoon}

Optical diffraction can blur these edges, motivating the blurred cartoon model
\eql{\label{eq-cartoon-model-smooth}
	\tilde\Th = \enscond{\tilde f = f \star h \in \Ldeux([0,1]^2)}{ f \in \Th \qandq h \in \Hh	}
}
where $\Th$ is the sharp-image model~\eqref{eq-cartoon-model} and $\Hh$ specifies the admissible blur kernels. For example, $h \geq 0$ may be required to be smooth and localized in both space and frequency.
Unknown blur makes it difficult to first detect the edges $\Gamma$ and then process the smooth regions separately.

Figure \ref{fig-sample-cartoon} shows examples of images in $\Th$ and $\tilde\Th$.

\myfigure{
\image{approximation}{.25}{sample-cartoon-1}
\image{approximation}{.25}{sample-cartoon-2}
\image{approximation}{.25}{sample-cartoon-3}
}{ 
	Examples of cartoon images: sharp discontinuities (left and center) and blurred edges (right).  
}{fig-sample-cartoon}




                                           
\section{Efficient Approximation}

                                        
\subsection{Decay of Approximation Error}
\label{sec-linear-approx-error}

To process signals in $\Theta$ efficiently, we seek an orthonormal basis in which the nonlinear approximation error $\norm{f-f_M}$ tends to $0$ as rapidly as possible when $M$ grows.


  
\paragraph{Polynomial error decay.}

This error decay is measured using a power law
\eql{\label{eq-approx-powerlaw}
	\foralls f \in \Th, \; \foralls M, \quad \norm{f-f_M}^2 \leq C_f M^{-\al}
}
Here $\al$ is common to all signals in the model; larger values mean faster decay. The exponent depends on the basis and the model $\Th$, and measures how efficiently the basis represents that class. The constant $C_f$ may depend on the particular signal $f$.

  
\paragraph{Relevance for compression, denoising and inverse problems.}

Approximation rates have practical consequences. Section \ref{sec-transform-coding} relates compression error to nonlinear approximation error, explaining why an effective approximation basis a